% ============================================================
% f2014_chongqing_science_q07_three_views.tex
% 2014 重庆卷（理科）第 7 题：三视图
% 来源 PDF：普通高考/2014/2014重庆理.pdf
% 原页：第 1 页（A3 横排）上栏，印面 865，第 7 题题干上方
% 说明：正视图为左高 5、右高 2、底长 4 的直角梯形；左视图为
%       宽 3、高 5 的矩形，内含一条从左上角斜向右下的虚线；
%       俯视图为两直角边长 4 与 3 的直角三角形。
% ============================================================
\begin{tikzpicture}[>=Stealth, line width=0.7pt, font=\small,
  x=1cm, y=1cm, line cap=round, line join=round,
  ext/.style={line width=0.4pt}]

  % ---------- 正视图 ----------
  \draw (0,0) -- (4,0) -- (4,2) -- (0,5) -- cycle;
  % 高 5
  \draw[ext] (0,0) -- (-0.6,0);
  \draw[ext] (0,5) -- (-0.6,5);
  \draw[<->] (-0.6,0) -- (-0.6,5) node[midway,left,inner sep=1pt] {$5$};
  % 底 4
  \draw[ext] (0,0) -- (0,-0.6);
  \draw[ext] (4,0) -- (4,-0.6);
  \draw[<->] (0,-0.6) -- (4,-0.6) node[midway,below,inner sep=2pt] {$4$};
  % 右高 2
  \draw[ext] (4,0) -- (4.6,0);
  \draw[ext] (4,2) -- (4.6,2);
  \draw[<->] (4.6,0) -- (4.6,2) node[midway,right,inner sep=1pt] {$2$};
  \node at (2,-1.5) {正视图};

  % ---------- 左视图 ----------
  \draw (6,0) rectangle (9,5);
  \draw[dashed] (6,5) -- (9,2);
  \draw[ext] (6,0) -- (6,-0.6);
  \draw[ext] (9,0) -- (9,-0.6);
  \draw[<->] (6,-0.6) -- (9,-0.6) node[midway,below,inner sep=2pt] {$3$};
  \node at (7.5,-1.5) {左视图};

  % ---------- 俯视图 ----------
  \draw (11,0) -- (15,0) -- (11,3) -- cycle;
  \node at (13,-1.5) {俯视图};
\end{tikzpicture}
